\documentclass[12pt]{article}

% === Packages ===
\usepackage[margin=1in]{geometry}
\usepackage{amsmath, amssymb}
\usepackage{setspace}
\usepackage{titlesec}
\usepackage{enumitem}
\usepackage{helvet}
\renewcommand{\familydefault}{\sfdefault} % modern sans-serif look

% === Formatting tweaks ===
\setstretch{1.15}
\setlength{\emergencystretch}{3em}
\setlist[enumerate]{itemsep=0.4em, topsep=0.2em}
\titleformat{\section}{\Large\bfseries\sffamily}{\thesection.}{1em}{}
\titleformat{\subsection}{\large\bfseries\sffamily}{\thesubsection}{1em}{}

% === Title ===
\title{\vspace{-2em}\textbf{Applied Mathematics Oral Exam Questions and Model Answers}\\
\large From Euler to Navier--Stokes}
\author{Prepared for Graduate Oral Examination}
\date{}

\begin{document}

\maketitle
\thispagestyle{empty}

% === Optional Intro Paragraph ===
These questions are designed for an applied mathematics oral examination following a research presentation on fluid dynamics.  
They emphasize conceptual understanding of the material derivative, incompressibility, the stress tensor, and the derivation and interpretation of the Navier--Stokes equations.  
Students are typically expected to demonstrate fluency in physical interpretation, modeling assumptions, and connections to research contexts rather than detailed on-board derivations.

Unless stated otherwise, bold symbols denote vectors or tensors, $\boldsymbol{\sigma}$ is the Cauchy stress, $\mathbf{b}$ is body force per unit mass, and the fluid is a classical single-constituent continuum with no couple stresses.

\section*{Applied Mathematics Oral Exam Questions and Model Answers:\\
From Euler to Navier--Stokes}

\subsection*{1. Kinematics and Material Derivative}
\begin{enumerate}
\item \textbf{Question:} Explain the difference between the local and convective components of the material derivative.\\
\textbf{Answer:} The local term $\partial \phi/\partial t$ measures change at a fixed spatial point. The convective term $(\mathbf{u}\cdot\nabla)\phi$ accounts for change caused by a parcel moving through spatial gradients. Together,
$D\phi/Dt=\partial_t\phi+\mathbf{u}\cdot\nabla\phi$ is the derivative following the motion. For a vector or tensor field, apply this operation componentwise.

\item \textbf{Question:} How does the material derivative help us interpret acceleration in a flow field?\\
\textbf{Answer:} The acceleration of a fluid particle is $D\mathbf{u}/Dt=\partial_t\mathbf{u}+(\mathbf{u}\cdot\nabla)\mathbf{u}$. Thus a steady Eulerian velocity field can still accelerate particles through spatial variation; ``steady'' does not mean zero particle acceleration.

\item \textbf{Question:} Give a physical example where the convective term dominates the local term.\\
\textbf{Answer:} In steady, spatially varying flow around an airfoil or through a nozzle, $\partial_t\mathbf{u}=\mathbf{0}$ while $(\mathbf{u}\cdot\nabla)\mathbf{u}$ can be large, so the acceleration is entirely convective.
\end{enumerate}

\subsection*{2. Incompressibility and Mass Conservation}
\begin{enumerate}
\item \textbf{Question:} What does $\nabla \cdot \mathbf{u} = 0$ mean physically?\\
\textbf{Answer:} It is the local condition for a volume-preserving velocity field: an infinitesimal material volume has zero instantaneous dilatation. It does not, by itself, say that the velocity is spatially uniform. Mass conservation is
$D\rho/Dt+\rho\nabla\cdot\mathbf{u}=0$; hence a positive-density material with $D\rho/Dt=0$ is divergence-free.

\item \textbf{Question:} How would you detect deviations from incompressibility?\\
\textbf{Answer:} Measure or compute both density variation and velocity divergence. In an incompressible numerical formulation, norms of $\nabla\cdot\mathbf{u}$ diagnose constraint error; in a compressible formulation, the continuity residual and $D\rho/Dt$ are also relevant. Measurement and differentiation errors must be distinguished from physical compressibility.

\item \textbf{Question:} Why is constant density often assumed even for slightly compressible fluids?\\
\textbf{Answer:} For many liquid flows and low-Mach gas flows with modest temperature and composition changes, relative density variations are small, so a constant-density model is accurate. Low Mach number alone is not sufficient in strongly heated, reacting, stratified, or variable-composition flows; an incompressible or low-Mach model must be justified by the density scales and physics.
\end{enumerate}

\subsection*{3. Pressure and Euler’s Equations}
\begin{enumerate}
\item \textbf{Question:} How does pressure arise as a surface traction?\\
\textbf{Answer:} With outward unit normal $\mathbf{n}$ and the convention $\boldsymbol{\sigma}=-p\mathbf{I}+\boldsymbol{\tau}$, the traction is $\mathbf{t}=\boldsymbol{\sigma}\mathbf{n}$ and its pressure part is $-p\mathbf{n}$. Thus $\int_{\partial V}(-p\mathbf{n})\,dS=-\int_V\nabla p\,dV$. The sign follows from the chosen stress convention.

\item \textbf{Question:} What assumptions distinguish Euler’s equations from Navier--Stokes?\\
\textbf{Answer:} Euler's equations assume an inviscid continuum with $\boldsymbol{\sigma}=-p\mathbf{I}$, whereas Navier--Stokes closes the stress with a viscous constitutive law. The familiar term $\mu\nabla^2\mathbf{u}$ follows only for a Newtonian fluid with constant viscosity and $\nabla\cdot\mathbf{u}=0$; more generally the viscous term is $\nabla\cdot\boldsymbol{\tau}$.

\item \textbf{Question:} Give a physical example where Euler’s equations are adequate and one where they fail.\\
\textbf{Answer:} Euler's equations can approximate the outer region of a high-Reynolds-number external flow away from viscous layers. They fail wherever viscous stress is leading order---for example, near a no-slip wall, in lubrication or creeping flow, and within dissipative shear layers. Even a thin viscous layer can determine global drag or separation, so ``high Reynolds number'' does not justify Euler everywhere.
\end{enumerate}

\subsection*{4. Cauchy Stress Tensor and Momentum Balance}
\begin{enumerate}
\item \textbf{Question:} What additional information does the Cauchy stress tensor contain compared to pressure?\\
\textbf{Answer:} The Cauchy stress maps a surface normal to traction, $\mathbf{t}(\mathbf{n})=\boldsymbol{\sigma}\mathbf{n}$. It contains direction-dependent normal and shear stresses; pressure is only its isotropic compressive part under the convention $\boldsymbol{\sigma}=-p\mathbf{I}+\boldsymbol{\tau}$.

\item \textbf{Question:} Why must the stress tensor be symmetric?\\
\textbf{Answer:} Local angular-momentum balance implies $\boldsymbol{\sigma}=\boldsymbol{\sigma}^{T}$ for a classical Cauchy continuum when there are no couple stresses or distributed body couples. Micropolar and other generalized continua need not have a symmetric Cauchy stress.

\item \textbf{Question:} Interpret each term in $\rho D\mathbf{u}/Dt = \nabla \cdot \boldsymbol{\sigma} + \rho \mathbf{b}$.\\
\textbf{Answer:} The left side is mass density times material acceleration, hence force per unit volume. The divergence $\nabla\cdot\boldsymbol{\sigma}$ is the net contact force per unit volume arising from tractions on an infinitesimal element, and $\rho\mathbf{b}$ is body force per unit volume, such as gravity when $\mathbf{b}=\mathbf{g}$.
\end{enumerate}

\subsection*{5. Viscous Stress and Newtonian Fluids}
\begin{enumerate}
\item \textbf{Question:} What is the physical meaning of the rate-of-strain tensor $\mathbf{D}$?\\
\textbf{Answer:} $\mathbf{D}=\tfrac12(\nabla\mathbf{u}+\nabla\mathbf{u}^{T})$ is the symmetric part of the velocity gradient. It measures instantaneous stretching and shear of material line elements; rigid-body rotation lies in the antisymmetric part and does not contribute to $\mathbf{D}$.

\item \textbf{Question:} What are the key assumptions behind $\boldsymbol{\tau} = 2\mu \mathbf{D} + \lambda (\nabla \cdot \mathbf{u})\mathbf{I}$?\\
\textbf{Answer:} This is the most general linear, isotropic, frame-indifferent viscous response of a simple Newtonian fluid, with instantaneous dependence on the local rate of deformation. The coefficients may depend on thermodynamic state but not on the deformation history. For incompressible flow the $\lambda$ term vanishes; for compressible flow thermodynamic admissibility constrains the shear and bulk viscosities.

\item \textbf{Question:} How does viscosity act as momentum diffusion?\\
\textbf{Answer:} Viscous stress transfers momentum across neighboring material layers. For constant $\mu$ in an incompressible Newtonian fluid, its divergence is $\mu\nabla^2\mathbf{u}$, a diffusion operator that damps velocity gradients and dissipates kinetic energy into internal energy.

\item \textbf{Question:} When does the Newtonian assumption break down?\\
\textbf{Answer:} It can break down in polymer solutions, concentrated suspensions, mucus, and blood at low shear rates, where stress may depend nonlinearly on shear rate, on deformation history, or on evolving microstructure. Whether a Newtonian approximation is acceptable depends on the material, time scale, and shear-rate range.
\end{enumerate}

\subsection*{6. Navier--Stokes Equations}
\begin{enumerate}
\item \textbf{Question:} Explain the physical meaning of each term in the incompressible Navier--Stokes equations.\\
\textbf{Answer:} In $\rho[\partial_t\mathbf{u}+(\mathbf{u}\cdot\nabla)\mathbf{u}]=-\nabla p+\mu\nabla^2\mathbf{u}+\rho\mathbf{b}$, the two left-hand terms are local and convective momentum change per unit volume. The right-hand terms are the pressure-gradient force, viscous momentum diffusion (under constant-$\mu$, incompressible assumptions), and body force. The companion constraint is $\nabla\cdot\mathbf{u}=0$.

\item \textbf{Question:} How does adding viscosity change the qualitative flow behavior?\\
\textbf{Answer:} Viscosity dissipates kinetic energy and diffuses momentum, regularizing velocity gradients over finite scales. Together with a separately specified wall condition, it produces viscous boundary layers. Viscosity does not itself ``enforce'' no-slip, and discontinuities in pressure or stress can still occur across interfaces carrying singular forces.

\item \textbf{Question:} In which regimes does the Reynolds number play a dominant role?\\
\textbf{Answer:} Reynolds number compares inertial and viscous scales for a stated choice of $U$ and $L$. Small $Re$ often supports a creeping-flow approximation; when all inertial terms are negligible, the steady Stokes equations are kinematically reversible. Large $Re$ permits strong scale separation and transition, but neither laminar/turbulent behavior nor steadiness is determined by $Re$ alone: geometry, forcing, stability, and disturbances matter.

\item \textbf{Question:} What does the term $\mu \nabla^2 \mathbf{u}$ represent physically?\\
\textbf{Answer:} It is the divergence of Newtonian viscous stress for constant $\mu$ and a divergence-free velocity. It diffuses each velocity component and dissipates kinetic energy. If viscosity varies or the flow is compressible, one must retain the more general expression $\nabla\cdot\boldsymbol{\tau}$.
\end{enumerate}

\subsection*{7. Boundary Conditions and Non-Dimensionalization}
\begin{enumerate}
\item \textbf{Question:} Explain the physical reasoning behind the no-slip and no-penetration conditions.\\
\textbf{Answer:} At an impermeable wall moving with velocity $\mathbf{u}_w$, no penetration is $(\mathbf{u}-\mathbf{u}_w)\cdot\mathbf{n}=0$. No slip additionally sets the tangential relative velocity to zero, so $\mathbf{u}=\mathbf{u}_w$. It is an effective interfacial model supported for many continuum liquid flows, not a consequence of incompressibility; slip laws may be appropriate at rarefied, microstructured, or specially treated interfaces.

\item \textbf{Question:} Why do we nondimensionalize the Navier--Stokes equations?\\
\textbf{Answer:} Nondimensionalization exposes competing balances, identifies independent dimensionless groups, permits similarity comparisons, and gives variables of order one for analysis and computation. It does not by itself guarantee numerical stability; the discretization and time integrator still determine stability.

\item \textbf{Question:} What insight does the Reynolds number provide?\\
\textbf{Answer:} With specified characteristic scales, $Re=UL/\nu$ measures the inertial-to-viscous scale ratio. It is an important control parameter and helps organize regimes, but critical values are configuration-dependent. Low $Re$ need not be steady, and high $Re$ need not be turbulent.
\end{enumerate}

\subsection*{8. Broader and Applied Questions}
\begin{enumerate}
\item \textbf{Question:} What numerical challenges arise at high Reynolds numbers?\\
\textbf{Answer:} High-$Re$ flows can contain thin shear layers and a broad range of dynamically active scales. A direct numerical simulation must resolve down to the dissipative scales and obey the time integrator's stability limits, so cost grows rapidly with $Re$. Under-resolution can cause aliasing and nonphysical energy accumulation; LES or RANS introduces modeling error instead. ``Stiffness'' depends on the formulation and solver rather than on $Re$ alone.

\item \textbf{Question:} How might the stress--strain relation change for complex or non-Newtonian fluids?\\
\textbf{Answer:} The stress may depend on strain history, rate magnitude, or additional internal variables, requiring constitutive laws such as power-law, Bingham, or viscoelastic models.

\item \textbf{Question:} In modeling biological or soft-matter flows, which assumptions must be revisited?\\
\textbf{Answer:} Revisit the constitutive law, wall or interface conditions, structural inertia, and coupling to elasticity, porosity, surface tension, chemistry, or active stresses. Liquid incompressibility is often still an excellent approximation, while effective viscosity or stress--strain relations may be non-Newtonian. The relevant assumptions should be checked against the system's scales rather than discarded automatically.

\item \textbf{Question:} How would you connect the continuum Navier--Stokes framework to molecular-scale descriptions?\\
\textbf{Answer:} Kinetic theory and nonequilibrium statistical mechanics define continuum fields as coarse-grained molecular moments and relate fluxes such as stress to microscopic dynamics. A continuum description also requires scale separation: the mean free path and molecular correlation scales must be small relative to the macroscopic length and time scales. Constitutive coefficients such as viscosity may be computed or measured from microscopic correlations.
\end{enumerate}
\end{document}
